% Ad Familiares 9.22 (ad-familiares-09-22) --- English translation
% Translated by Claude (Anthropic) from the Latin (Perseus).
% Style guide: STYLE.md.

\ciceroLetterOpener{CICERO PAETO.}

\ciceroSection{1} I love your modesty --- or rather, your freedom of speech. And yet this was the view of Zeno, a sharp fellow, by Hercules, even if our Academy has a great quarrel with him --- but, as I say, the Stoics hold that each thing is to be called by its own name. For their reasoning runs like this: nothing is obscene, nothing shameful to utter; for if there is any indecency in obscenity, it lies either in the thing or in the word; there is no third option. In the thing it does not lie. And so in comedies the act itself is narrated (as in that scene in \textit{The Craftsman}: \textit{just now, by chance ---} \,you know the song; you remember Roscius: \textit{so she left me high and dry, naked ---}: the whole speech is veiled in words, the matter being more brazen than what is said); and so too in tragedies. For what is that line, \textit{the one woman who ---}? what, I ask, of \textit{she usurps the double couch}? what of \textit{this fellow's \dots dared to enter his bed}? what of \textit{when I was still a virgin, against my will, by violence Jupiter ravishes me}? Good, ``ravishes''; though it means the very same thing, no one would have stood the other word.

\ciceroSection{2} You see, then, that since the matter is the same and only the words differ, nothing seems shameful. So: it is not in the thing; much less is it in the words. For if what the word signifies is not shameful, the word that signifies it cannot be shameful. You say ``old woman'' (\textit{anus}) by a borrowed name; why not by its proper one? If the thing is shameful, not even the borrowed name is clean; if it is not, then the proper one is better. The ancients used to call the tail (\textit{cauda}) ``penis'' --- whence, from the resemblance, ``penicillus,'' a little brush; but today ``penis'' belongs to the obscenities. ``But Piso Frugi, in his \textit{Annals}, complains that young men are `given over to the \textit{penis}.''' What you in your letter called by its own name he, more veiled, called \textit{penis}; but, because many followed suit, the word has become as obscene as the one you used. What about that common expression, \textit{cum nos te voluimus convenire} --- ``when we wanted to meet you'' --- is that obscene? I remember an eloquent ex-consul speaking thus in the Senate: \textit{hanc culpam maiorem an illam dicam?} --- ``shall I call this fault greater or that?'' Could anything sound more obscene? ``No,'' you say; ``for he did not mean it that way.'' So it is not in the word; and I have shown it is not in the thing; therefore it is nowhere.

\ciceroSection{3} \textit{Liberis dare operam} --- ``to devote oneself to (begetting) children'' --- how honourably it is said! Even fathers ask their sons to do it; the name of the act itself they dare not utter. Socrates was taught the lyre by the most distinguished lyre-player; his name was Connus --- you don't think \emph{that's} obscene, do you? When we say \textit{terni}, ``three apiece,'' we are saying nothing scandalous; but when we say \textit{bini}, ``two apiece,'' it is obscene. ``In Greek, anyway,'' you will say. So it is not in the word, since I too know Greek and yet I say \textit{bini} to you, and you do likewise, just as if I had spoken Greek and not Latin. \textit{Ruta} (rue) and \textit{menta} (mint) are both proper words; I would like to call a little mint-plant \textit{mentula}, on the model of \textit{rutula} --- I cannot. ``\textit{Tectoriola}'' --- little plasters --- is charming; on that pattern, then, say ``\textit{pavimenta}'' too --- you can't. Do you see, then, that there is nothing here but silliness, that shamefulness lies neither in the word nor in the thing, and so is nowhere?

\ciceroSection{4} And so we plant obscenities inside respectable words. How so? Is \textit{divisio}, ``division,'' not a respectable word? Yet there is an obscenity inside it; and the same with its counterpart \textit{intercapedo}, ``interruption.'' Are these words obscene, then? But we behave absurdly: if we say ``so-and-so strangled his father,'' we make no apology; but if it is something about Aurelia or Lollia, an apology must be premised. And by now even non-obscene words are taken for obscene. ``\textit{Batuit}, `he gave it a beating,' '' he says, ``is shameless; \textit{depsit}, `he kneaded,' is much more shameless.'' And yet neither is obscene. The world is full of fools. ``\textit{Testes}'' is a perfectly respectable word in a court of law; in another setting, less so. And ``the \textit{colei Lanuvini},'' the Lanuvine balls (of cheese), are respectable; the Cliternine ones, not so much. What of this --- the thing itself is sometimes respectable, sometimes not. \textit{Suppedit}, ``he pisses underneath himself,'' is disgraceful; but the moment he is naked in the bath, you will not find fault.

\ciceroSection{5} There you have the Stoic doctrine: \textit{the wise man will call things by their straight names} \textit{[Greek: ho sophos euthyrr\=emon\=esei]}. What a lot has come out of one word of yours! Your daring me in every direction is welcome; for my part, I keep, and shall keep (such is my habit), Plato's reserve. So I have written to you in veiled words about things the Stoics handle with the plainest words possible --- though they say that breakings of wind ought to be just as free as belchings. So an honour, then, to the Kalends of March! You will love me and keep well.
